\documentclass[12pt,letterpaper]{article}
\usepackage[utf8]{inputenc}
\usepackage[spanish]{babel}
\usepackage{amsmath}
\usepackage{amsfonts}
\usepackage{amssymb}
\usepackage[margin=2.5cm]{geometry}

\begin{document}

\begin{center}
    \Large \textbf{Evaluación de Examen 2: Unidad II} \\
    \large Física para Ciencias de la Salud (005-1713) \\
    \vspace{0.5cm}
    \textbf{Estudiante:} Johannys Rojas \quad \textbf{C.I.:} 32.828.239
    \vspace{0.3cm} \\
    \large \textbf{Calificación Final: 10/10 puntos}
\end{center}

\vspace{0.5cm}

\section*{Análisis de Resultados}

\subsection*{1. Cinemática (Aceleración media)}
\textbf{Procedimiento y resultado:} Extrae los datos correctamente y plantea la ecuación de aceleración media. Sustituye y efectúa la división de $0,6 \text{ m/s} / 0,15 \text{ s}$, obteniendo el resultado analítico exacto de $4 \text{ m/s}^2$ con un correcto manejo de las unidades. \\
\textbf{Veredicto:} \textbf{Correcto} (2/2 pts).

\subsection*{2. Dinámica (Leyes de Newton)}
\textbf{Procedimiento y resultado:} Muestra un despeje impecable a partir de la Segunda Ley de Newton ($a = \frac{F}{m}$). Sustituye los valores y divide $150 \text{ N}$ entre $25 \text{ kg}$, logrando un resultado perfecto de $6 \text{ m/s}^2$. \\
\textbf{Veredicto:} \textbf{Correcto} (2/2 pts).

\subsection*{3. Trabajo Mecánico}
\textbf{Procedimiento y resultado:} Justifica de gran manera el uso de la fórmula general del trabajo, indicando de antemano explícitamente que $\cos(0^\circ) = 1$. La multiplicación ($400 \text{ N} \cdot 0,8 \text{ m}$) es perfecta y su resultado de $320 \text{ Joules}$ es completamente exacto. \\
\textbf{Veredicto:} \textbf{Correcto} (2/2 pts).

\subsection*{4. Energía Potencial}
\textbf{Procedimiento y resultado:} Extrae de manera estructurada los datos del enunciado y aplica a la perfección la fórmula de la energía potencial gravitatoria ($E_p = m \cdot g \cdot h$). Destaca muy positivamente que desarrolla un paso intermedio calculando el peso en Newtons ($11,76 \text{ N}$) antes de multiplicarlo por la altura ($1,5 \text{ m}$), lo cual evidencia un sólido entendimiento analítico dimensional. El resultado final es $17,64 \text{ Joules}$. \\
\textbf{Veredicto:} \textbf{Correcto} (2/2 pts).

\subsection*{5. Potencia Mecánica}
\textbf{Procedimiento y resultado:} Plantea correctamente la fórmula de potencia mecánica y relaciona el trabajo entre el tiempo ($P = \frac{W}{t}$). Ejecuta la sustitución de $75 \text{ J} / 60 \text{ s}$ logrando el resultado perfecto de $1,25 \text{ W}$. \\
\textbf{Veredicto:} \textbf{Correcto} (2/2 pts).

\vspace{0.5cm}
\section*{Comentarios Generales y Sugerencias}
¡Muchas felicidades por un examen perfecto! La estudiante demuestra un entendimiento profundo tanto de los conceptos físicos de la Unidad II como de los procedimientos matemáticos necesarios.
\begin{itemize}
    \item El orden al separar sistemáticamente los \textit{Datos} y la \textit{Fórmula} antes de operar permite seguir el desarrollo analítico sin dificultades, lo cual es el estándar ideal en la documentación científica y de la salud.
    \item Es destacable y muy meritorio el excelente manejo del \textbf{análisis dimensional}, especialmente notable en la pregunta 4 al plasmar de manera explícita la obtención de los Newtons como paso previo antes de llegar a los Joules.
    \item Se incentiva a mantener este fantástico nivel de organización, limpieza y precisión en las evaluaciones venideras.
\end{itemize}

\end{document}